% !TEX root = ../b-rg.tex
\chapter{Word Formation}\label{ch:word-formation}
\stub

% --- Planned sections ---
% Organised by the class of the new word, then by the class it is formed
% from (Lingua questionnaire 2.2).
% \section{Nouns}
%   From verbs: -ment (porçonment), -çon, -enç (from -ir verbs), agent
%   -our (forniscour), conversion (desir). From adjectives: -tað / -stað
%   (-ous -> -ostað). From nouns: collective -ary (volary), diminutives.
% \section{Adjectives}
%   -er (marcer), -ous, -abr '-able' (also the potential form, cf. TAM),
%   -esc (Borallesc, moresc); participles.
% \section{Verbs}
%   Denominal -ar verbs (porçon -> porçonnar); prefixes des- (desjeun ->
%   desjeunar), re-, in-, entr- (entrcollocq), sur- (surcrescr) ...
% \section{Adverbs}
%   cos + adjective (cos nuvr 'lately'); other sources.
% \section{Compounding}
%   Noun + adjective and noun + de + noun lexicalised phrases
%   (sar marcer, tort dell'ivan); hyphenated compounds (doutr-alp);
%   complex prepositions (cf. Prepositions).
% \section{Inherited, learned and borrowed vocabulary}
%   Doublets (inherited vs learned); Celtic, Old English, Norse, Vascon,
%   Arabic strata; how loans are nativised in spelling.
